%\pagebreak
\section{Self Mapping}
\label{sec:self_mapping}

\index{self maps}

In general, a map operation may create a separate copy of a list item in a
device data environment that the runtime implementation must associate with the
original list item. A \emph{self map} is a map operation that does not result in a
separate copy of a list item. Rather, the original storage of the list item is
reused for the device data environment, with the runtime regarding its storage
as being ``mapped'' to itself.

OpenMP 6.0 adds the \kcode{self} modifier to the \kcode{map} clause to require
that its list items are mapped using a self map.  The \kcode{self} keyword was
also added as an argument to the \kcode{defaultmap} clause, indicating that
variables of a specified category in a \kcode{target} construct are self mapped
by default. Additionally, the \kcode{self_maps} clause was added to the
\kcode{requires} directive to require that all map operations in a given
compilation unit should be self maps. The \kcode{self_maps} clause includes all
the guarantees provided by the \kcode{unified_shared_memory} requirement clause.

The following C example shows the use of these features when mapping a structure
type to a device. The self map allows \ucode{start} and \ucode{end} pointer
data members that point to the start and end of the \ucode{buf} data member to
be used on the device without requiring pointer attachments.

\cexample[6.0]{self_map}{1}

The next C++ and Fortran example maps a structure for which an array member
\ucode{buf} is private but may be accessed in the public interface via a
pointer member \ucode{p}. When mapping the structure to a device, the
programmer must ordinarily ensure that the \ucode{p} data member on the device
is attached to the \ucode{buf} data member, by adding a \ucode{my_data.p[:]}
(for C++) or \ucode{my_data\%p(:)} (for Fortran) list item to a \kcode{map}
clause. By instead asking for the structure to be self mapped, there is no need
for the pointer attachment.

\cppexample[6.0]{self_map}{2}
\ffreeexample[6.0]{self_map}{2}

If the implementation is unable to satisfy the requirements of a self map then
a runtime error will be issued. This may occur because the original storage is
not accessible and cannot be made accessible from the device, or it may occur
because the storage has already been mapped to the device without a self map. A
third case is that the self map applies to a pointer for which pointer
attachment is prescribed to a pointee that was not also self mapped. Since a
self-mapped attached pointer in that case would assign a device address to the
original pointer which would almost certainly not be the desired behavior.

The following example illustrates each of the three cases above that could
potentially result in a runtime error due to an unfulfilled self map.

\cexample[6.0]{self_map}{3}
\ffreeexample[6.0]{self_map}{3}

